\SourceLeaf{339}%
ergeben. Auf \(\log\tau\) hat das keinen Einfluss, sodass wir \(\log t=2{,}6238578,\ t=420''\ 5889,\ T=5^0 38'51''4755\) erhalten. Würde man die Rechnung mit diesem Wert von \(T\) noch einmal durchführen, so bliebe \(B\) unverändert. Für \(\log\tau\) erhielte man \(2{,}6238470\), also \(\log t=2{,}6238573,\ t=420''\ 5884\) und damit \(T=5^0 38'51''4757\). Eine weitere Rechnung mit diesem Wert würde gar keine Änderung mehr ergeben. Offenbar hätte man bereits beim vorigen Ergebnis stehen bleiben können, da bei siebenstelligen Logarithmen die vierte Dezimalstelle der Sekunde noch um eine oder einige Einheiten schwanken kann. Das Endergebnis lautet also
\[
\begin{gathered}
\text{Breite von Inselsberg}=B^0-b=50^0\ 51'\ 8''\ 9444\\
\text{Azimut der Seite Inselsberg-Brocken}=180^0+T^0-t\\
=185^0\ 35'\ 21''\ 1815
\end{gathered}
\]
Für den Längenunterschied findet man schließlich
\[
\begin{aligned}
\log\lambda&=2{,}7313519\\
\log l&=2{,}7313472\\
l&=538''7002=0^0\ 8'\ 58''\ 7002
\end{aligned}
\]

Die Bemerkungen am Schluss von Artikel 20 brauchen hier nicht wiederholt zu werden, da sie auch hier ihre volle Gültigkeit behalten.
\par\medskip\centerline{\rule{.25\textwidth}{.3pt}}
